\section{A finite cutoff for unrestricted witnesses}
Let $S=i+k+\ell$, $w=a^ibb a^kb a^\ell$, and assume $k>0$ and
$i+\ell>0$. Write $d=x-u>0$, $D=v-y=Sd/3$. We may orient an
unrestricted witness this way by exchange; under Holub's condition the
length identity already gives this orientation for distinct substitutions.
Indeed $d=0$ gives identical images by reading corresponding occurrences.

\begin{theorem}
Define the following explicit integers:
\[
 n=1000(S+10),\qquad q=1000n,\qquad
 N=(100q)^{4q},\qquad d_0=(204n)^N .
\]
If there is an unrestricted counterexample, there is one with
$0<d\le d_0$. For $k>2(i+\ell)$, the same implication
holds with Holub's condition on both the premise and the conclusion.
\end{theorem}

\begin{proof}
First choose a noncommuting ordered pair $E,F$ of the four images.
The equal-length words $EF,FE$ disagree at some letter. A letterwise
projection to a binary alphabet preserving that mismatch preserves all
lengths and the equation. Choose constants $b_0,b_1$ for the two
different letters at the first mismatch (either order is allowed).
Add to the original equation the equations
\[
 EF=P b_0 Q,\qquad FE=P b_1 T .
\]
If $i>0$, $U$ is a proper prefix of $X$, so add $X=UcZ$;
if $i=0$ use $X=ZcU$ since $\ell>0$. Here $c$ is one of the two
constants. Conversely each branch gives $d>0$ and a genuine mismatch.

For clarity we give numerical constants in the small-solution bound.
Combine the four equations in two rounds, using for a pair
$L=R$, $L'=R'$ the single equation
\[
 L b_0 L' L b_1 L'=R b_0 R' R b_1 R'.
\]
Both halves on either side have equal length, hence the left and right
half lengths agree. Comparing the two split equalities, the unique
position at which the halves differ is $|L|$ on the left and $|R|$
on the right. Thus $L=R$ and $L'=R'$ by cancellation. The original
sum of side lengths is $2(S+3)+5+5+4$. Each round doubles this
sum and adds four per combination. Including the five constant
symbols of the involutive alphabet below, the input size is at most $n$.

Apply Ciobanu--Diekert--Elder, \emph{Solution sets for equations over
free groups are EDT0L languages}, arXiv:1508.02149, Theorem~4,
in its free \emph{monoid} case. Use the two positive binary constants,
their formal inverses and the marker. Any positive witness is reduced.
Any returned solution in reduced words solves literal monoid equalities;
mapping each inverse letter to its positive counterpart preserves the two
marked constants, every length, and every branch equation.

Here are details of the state count. In that construction, Definition~8
and Section~3.10.4 give state data $(W,B,\mathcal X,\theta,\mu)$:
$|W|\le204n$, with at most $615n$ constants and $6n$ variables.
There are fewer than $(q+1)^{204n+1}$ choices of word, at most $2^q$
choices of the active symbols (the constant and variable pools are
disjoint), at most $(q+1)^q$ partial type maps, and at most $19^q$
partial constraint maps. Indeed the constraint monoid consists of the
sixteen pairs of signed binary letters, zero, and the identity.
Their product is below $N$. Section~3.7 labels substitution edges by the
identity on constants. Compression edges expand a letter by at most two
except on the final transition, where the sum of active image lengths
is at most $|W|\le204n$; inactive symbols are fixed.
A nonempty accepting path can be shortened by deleting cycles, giving
fewer than $N$ edges. The theorem reads each variable image by composing
these labels on a single distinguished letter. Every resulting length
is at most $(204n)^N$. All branch equations still hold on this path,
so this bounds $d$ as claimed. The constrained assertion follows from
the middle-block lemma and the edge-case argument below.
\end{proof}

\begin{lemma}
Let $k>2(i+\ell)$, $i+\ell>0$, $u>0$.
If $ku\ge x+u-\gcd(x,u)$, every solution with $d>0$ commutes.
\end{lemma}
\begin{proof}
In the middle $X^k$, the subword $U^k$ begins after $L=2D-id$
letters and ends $Q=D-\ell d$ letters before the end. Both numbers
are positive by hypothesis. The factor of length $ku$ has periods
$x,u$, so Fine--Wilf gives period $g=\gcd(x,u)$. Since it contains
$x$ and $u$ consecutive letters, the bi-infinite periodic extensions
of the two blocks agree with a common $g$-periodic sequence there,
possibly in different phases at their respective block starts.
Consequently their primitive roots are rotations of one another of
length $r\mid g$. The outer run with positive exponent makes $U$ a
prefix or suffix of $X$ (as $u<x$). As $r$ divides both lengths, this
aligns their roots: $X,U\in R^*$ for one primitive $R$.
The occurrence at offset $L$ now starts with a full $R$, whence $r\mid L$.
Since $r\mid d$, this gives $r\mid 2D$, and $3D=Sd$ gives $r\mid D$.
In particular $r\mid Q$. Read from the end of $U^k$ to the end of
the equation and cancel $U^\ell$. This gives
\[
 V=R^{Q/r}Y R^{\ell d/r}.
\]
If $R,Y$ do not commute, they constitute a two-word code (the standard
two-word defect theorem). After the substitution, the left side has
adjacent $Y$ tokens and the first two right-side $Y$ tokens are
separated by $D/r>0$ copies of $R$. This contradicts unique decoding.
Thus $Y$, and then $V$, commute with $R$.
Finally if $k\ge4$ and $v<x$, then $u>(S/3-1)d\ge2d/3$, since
$S\ge5$. Hence $(k-2)u>d$, implying the required Fine--Wilf bound.
It follows that every oriented unrestricted counterexample in this
subfamily already satisfies $v\ge x>u$ and $v\ge y$.
\end{proof}

\paragraph{The two edge tuples.}
For $k>2(i+\ell)$ with $k<4$ the only possibilities have $k=3$ and
$i+\ell=1$. Write $d=3s$, $D=4s$ with $s>0$.
We show that in every counterexample a full copy of $X$ occurs in $V^2$.
Indeed the first $L=2D-id$ letters of the middle $X^3$ are in $V^2$
by the start and end positions above. If $i=0$ and $x>L=8s$ then
$u>5s>d$, so the middle-block lemma applies.

Suppose $i=1$, $\ell=0$, $x>L=5s$. Write $X=UT$, $|T|=3s$.
The initial $V$ starts with $T$ because it starts just after the initial
$U$ and $v\ge4s$. The final $V$ is $AY$ where $A$ is the suffix of
$X^3$ of length $4s$. Since $x>5s$, $A$ is a suffix of $X$.
Thus $A$ begins and ends with $T$, so it has period $s$ and equals
$C^4$ for a word $C$ of length $s$. In particular $X=UC^3$.
The middle-block lemma covers $u\ge3s$, so assume $2s<u<3s$.
From the nested middle occurrence, $U^3 C^4$ is a suffix of $X^3$.
Cancel $C^3$ at the right: $U^3 C$ is a suffix of $X^2 U$.
Consequently $U$ is a suffix of $U^3 C$, giving period $s$ on $U$
and suffix $C$. Inside the middle $X^3$, the second $X$ starts
$x-L=h=u-2s$ letters after the start of $U^3$. It starts with $U$,
and $0<h<s$, so the length-$u$ window in $U^3$ at offset $h$ equals
$U$. The primitive root length $r$ of $U$ divides $h$, by invariance
under this rotation. Fine--Wilf on $U$ with periods $s,r$ now gives
$r\mid s$ since $u>2s>s+r$. Hence $C$, its length-$s$ suffix, has
the same primitive root, and $X,U,A$ are powers of it. As $V=AY$ with
$A$ a nonempty such power, the adjacent-$Y$ code argument above forces
$Y$ to commute as well. This proves the full-copy claim.

For a Holub-admissible counterexample, $x\le v$, so a copy of $X$
in $V^2$ is a prefix of one rotation of $V$. Add the two equations
$V=JK$, $KJ=XH$ to the branch. These equations themselves guarantee
$v\ge x$ on every returned solution. Conversely the copy just proved
supplies $J,K,H$ in every admissible counterexample. Pad these six
equations with two empty identities and combine in three rounds.
The original total side length is at most $2S+27$; the successive
totals are at most $2(2S+27)+16$, twice that plus eight, and twice
the latter plus four. Including the alphabet still gives input size
below $n$. Pad the system
to eight equations with empty equalities and use three combination
rounds. The input is still below $1000(S+10)$: the two new equations
add seven symbols, each original symbol is copied eight times, and
all combination constants together contribute $48$. The same numerical
bound therefore applies.

\paragraph{A smaller exact box at fixed $d$.}
For one input letter let $A,B$ be its left and right images, of lengths
$a,b$, and $m,M$ its extreme start displacements (left minus right).
Write $P=M-m>0$, $z=\max(0,m)$, $t=z-m$.
A minimum unrestricted counterexample cannot satisfy all three bounds
\[
 a\ge t+2P,\qquad a\ge t+P+d,\qquad b\ge z+2P+d.
\]
Indeed $B[j]=A[j-m]=B[j+P]$ throughout $[z,z+P+d)$.
Thus $B[z:z+2P+d]$ has period $P$. Cut $P$ letters at $t$ in
every left copy and at $t+e$ in each corresponding right copy of
displacement $e$. These are paired cuts in the common word. Their
starts in $B$ range from $z$ to $z+P$, and all yield the same shortened
$B$ because deleting one period anywhere between these starts gives
the same word. The deleted block at $z$ can be recovered from
$B[z+P:z+2P]$. In the left image the two adjacent $P$-blocks at $t$
are equal as well: the occurrence with displacement $m$ copies them
from $B$, using $a\ge t+2P$. If the shortened solution commutes,
its primitive root has length $r\mid d$ since $d$ was preserved.
The retained $P+d$ letters in the right image have periods $P,r$.
Fine--Wilf and the presence of a full primitive rotation give $r\mid P$.
Inserting the following $P$ letters again in either image thus preserves
the periodic phase. This contradicts the original noncommutation.
The old and new bounds can be intersected since each applies to the
same minimum unrestricted witness.

\paragraph{Exact component criterion.}
At any fixed lengths let $h=\gcd(x,y,u,v)$. Every position equality
preserves the offset modulo $h$ inside the corresponding image; upon
division by $h$ the graph is $h$ interleaved copies of the normalized
graph, with identical adjacencies. If the normalized graph is connected,
every image is a power of the same length-$h$ word. Conversely if it
is disconnected assign distinct symbols to normalized components.
Commuting images would have a primitive root whose length divides all
four normalized lengths, hence would be unary, impossible. Replacing
each symbol by $h$ copies lifts this witness injectively. This proves
both the test and a genuine $d\mapsto d/h$ descent for $h>1$, preserving
Holub's inequality. Literal equality and inequality of words suffice to
verify each resulting witness.

\paragraph{Finite verification.}
For all 210 residual tuples through exponent six, the exact computation
in the accompanying checks excludes every unrestricted counterexample
with $1\le d\le12$. In particular it excludes every admissible one.
For each integrally possible $D$ it uses $x=u+d,v=y+D$ and enumerates
the intersection of the two minimum-witness boxes, independently in $u,y$.
Impossibility of an unrestricted minimum excludes the entire original box
and all lengths beyond it at that same $d$.

\paragraph{Endpoint order alone loses information.}
For the nonforcing control $w=abbaba$, the lengths $(x,y,u,v)$
given by $(2,1,1,2)$ and $(3,1,1,3)$ both satisfy Holub's inequality
and have gcd one. Their cumulative left and right endpoints are,
respectively, $(2,3,4,6,7,9),(1,3,5,6,8,9)$ and
$(3,4,5,8,9,12),(1,4,7,8,11,12)$.
All cross-comparisons have the same sign. In the first case write
$U=p,Y=q$. Reading the first two letters gives $X=pV[0]$;
the third letter gives $V[1]=q$, the fourth gives $V[0]=q$,
and the fifth gives $p=q$. All images are unary.
In the second case $(X,Y,U,V)=(101,0,1,010)$ is an admissible
noncommuting solution by direct concatenation. Thus even the addition
of the maximum-length test and primitive length normalization does not
make counterexample existence constant on the endpoint-order cells.

\paragraph{Coverage check.}
A second run directly joins positions for every length vector in the
original unrestricted rectangle $0\le u<T_a$, $0\le y<T_b$.
For the admissible run, keep the rows with $y+D\ge u+d$, and append
all such rows with
$T_b\le y<\max(T_b,T_a+d+3P_b-D)$ at the same $u$.
This is precisely the previous section's admissible box.
In each mode the minimum-counterexample argument removes the upper
length restriction at fixed $d$. The exponent-four run is also tabulated
independently of the exponent-six run.
